%=====================================================================
\chapter{Glossary}\label{app:glossary}
%=====================================================================
\normalfigures

Every technical term introduced in this handout, with the chapter that defines it.

\newcommand{\gt}[3]{\textbf{#1}~{\small\textit{(Ch.~#2)}} --- #3\par\smallskip}

\begin{multicols}{2}
\small
\raggedright

\gt{A/B test}{10}{A randomised experiment comparing two variants on live traffic;
the practical form of a randomised controlled trial.}

\gt{Accuracy}{14}{Proportion of all predictions that are correct. Misleading
whenever classes are imbalanced.}

\gt{Activation function}{17}{The non-linear function applied to a neuron's weighted
sum. Without it, depth gains nothing.}

\gt{Adversarial machine learning}{26}{Attacks against an ML system itself: evasion,
poisoning, model extraction, membership inference.}

\gt{Agent}{22}{A system that pursues a goal over multiple steps, choosing its own
actions and calling tools.}

\gt{Alert fatigue}{26}{Degradation of response quality when alert volume exceeds
what analysts can investigate.}

\gt{Anomaly detection}{16}{Identifying observations that differ markedly from the
majority; used when the rare class has no labels.}

\gt{Attention}{18}{A mechanism letting every position in a sequence consult every
other directly, with learned weights.}

\gt{Autoencoder}{18}{A network trained to reconstruct its input through a
bottleneck; reconstruction error serves as an anomaly score.}

\gt{Bagging}{15}{Training many models on bootstrap samples and averaging them.
Reduces variance.}

\gt{Base rate}{6}{The underlying prevalence of a class. Low base rates make even
excellent detectors imprecise.}

\gt{Baseline}{11}{The simplest sensible rule. A model that does not beat it is not
a result.}

\gt{Batch}{17}{The group of examples processed before one weight update.}

\gt{Bayes' theorem}{6}{Converts a measurable likelihood into the posterior you
actually want.}

\gt{Bias (statistical)}{11}{Error from a model too rigid to represent the truth.}

\gt{Bias (societal)}{28}{Systematic unfairness towards a group, entering at any of
five stages from the world to deployment.}

\gt{Boosting}{15}{Training models sequentially, each fitted to the previous one's
residuals. Reduces bias.}

\gt{Bottleneck}{18}{The narrow layer of an autoencoder that forces compression.}

\gt{Calibration}{28}{A score of 0.7 corresponds to a 70\% observed rate --- and does
so equally in every group.}

\gt{Canary release}{27}{Serving a new model to a small share of traffic with
automatic rollback.}

\gt{Central Limit Theorem}{6}{Sample means approach a normal distribution
regardless of the population's shape; spread falls as one over root $n$.}

\gt{Classification}{13}{Supervised prediction of a discrete label.}

\gt{Clustering}{16}{Unsupervised grouping of similar observations.}

\gt{Collider}{10}{A variable caused by both X and Y. Controlling for it
\emph{creates} a spurious association.}

\gt{Concept drift}{19}{The relationship between features and target changes over
time, degrading a deployed model.}

\gt{Confidence interval}{9}{A range constructed so that the procedure captures the
true value a stated proportion of the time.}

\gt{Confounder}{10}{A variable causing both X and Y, creating association without
causation. Control for it.}

\gt{Confusion matrix}{14}{The table of TP, FP, FN, TN from which all classification
metrics derive.}

\gt{Container}{24}{An application packaged with its libraries, sharing the host
kernel. The practical basis of reproducibility.}

\gt{Convolution}{18}{Sliding a small learned kernel across a grid input, sharing
weights across all positions.}

\gt{Correlation}{8}{Strength and direction of a \emph{linear} association.}

\gt{Cosine similarity}{4}{Similarity of direction, ignoring magnitude. The standard
measure for text.}

\gt{Cross-validation}{14}{Repeated train/validate splits so every observation is
validated once.}

\gt{CRISP-DM}{2}{The six-phase cyclic lifecycle from business understanding to
deployment.}

\gt{Curse of dimensionality}{15}{As features multiply, data becomes sparse and all
points nearly equidistant, breaking distance-based methods.}

\gt{DAG}{23}{Directed acyclic graph; the dependency structure of a pipeline.}

\gt{Dashboard}{31}{A decision-support display. Design it around a decision, not
around available data.}

\gt{Data contract}{23}{A versioned agreement on schema, ranges and freshness,
enforced automatically at ingestion.}

\gt{Data lake}{23}{Storage for raw data of any type, with schema applied at read
time.}

\gt{Data warehouse}{23}{Storage for cleaned, modelled tables, with schema fixed at
write time.}

\gt{DBSCAN}{16}{Density-based clustering that finds arbitrary shapes and labels
sparse points as noise.}

\gt{Decision boundary}{13}{The surface at which a classifier's predicted class
changes.}

\gt{Decision tree}{13}{A model that splits data by a sequence of feature tests
chosen to reduce impurity.}

\gt{Deep learning}{17}{Machine learning with multi-layer neural networks that learn
their own features.}

\gt{Demographic parity}{28}{Equal flag rates across groups. Incompatible with
calibration when base rates differ.}

\gt{Digital twin}{25}{A live software model of a physical asset, updated from its
telemetry.}

\gt{Dropout}{17}{Randomly silencing neurons during training to force redundancy.}

\gt{Edge computing}{24}{Computation on or beside the device producing the data.
Lowest latency, tightest resources.}

\gt{Effect size}{9}{How large a difference is, independent of sample size. Report
with every $p$-value.}

\gt{Elbow method}{16}{Choosing the cluster count by locating the bend in the
inertia curve.}

\gt{Embedding}{20}{A dense vector representation in which semantically similar
items are geometrically close.}

\gt{Ensemble}{15}{Many models combined; works because their errors differ.}

\gt{Epoch}{17}{One complete pass over the training data.}

\gt{Equal opportunity}{28}{Equal recall across groups: those who need help have an
equal chance of being identified.}

\gt{ETL / ELT}{23}{Extract-Transform-Load versus Extract-Load-Transform. ELT
dominates because storage is cheap.}

\gt{Evasion}{26}{Perturbing an input at inference time to escape detection.}

\gt{Explainability}{28}{The ability to say why a model produced a particular
output, in terms the audience can use.}

\gt{F1 score}{14}{Harmonic mean of precision and recall; appropriate when both
errors cost similarly.}

\gt{Feature}{3}{One measured attribute; a column.}

\gt{Feature engineering}{15}{Constructing informative features from raw data.
Usually worth more than changing the algorithm.}

\gt{Feature store}{27}{A single service computing features for both training and
serving, eliminating skew.}

\gt{Fog computing}{24}{An intermediate tier --- a gateway serving one site.}

\gt{Generative AI}{21}{Systems that create new content by predicting plausible
continuations.}

\gt{Gini impurity}{13}{A node purity measure used to choose tree splits; zero for a
pure node.}

\gt{Gradient}{4}{The vector of partial derivatives; points in the direction of
steepest increase.}

\gt{Gradient descent}{4}{Repeatedly stepping against the gradient; the optimisation
underlying almost every model here.}

\gt{Ground-truth lag}{27}{The delay before the true outcome is known. Drives the
design of monitoring.}

\gt{Hallucination}{21}{Fluent, confident, fabricated output --- a consequence of
next-token generation, not a bug.}

\gt{Heteroscedasticity}{12}{Error variance that changes with the prediction;
visible as a fan in the residual plot.}

\gt{Hyperparameter}{11}{A setting chosen by you rather than learned.}

\gt{Idempotence}{23}{Re-running a task produces the same result rather than
duplicating its effect.}

\gt{Imputation}{7}{Filling missing values. Inappropriate under MNAR, where absence
is itself informative.}

\gt{Inertia}{16}{Total squared distance from points to their assigned centroid.}

\gt{Isolation forest}{16}{Anomaly detection by random splitting; the strong default
when no labels exist.}

\gt{Interval scale}{3}{Numeric with equal gaps but an arbitrary zero; differences
are meaningful, ratios are not.}

\gt{k-means}{16}{Clustering by alternating assignment to nearest centroid and
recomputation of centroids.}

\gt{k-nearest neighbours}{13}{Classification by majority vote among the closest
training points. Requires scaling.}

\gt{Lasso (L1)}{12}{Regularisation that drives coefficients exactly to zero,
performing feature selection.}

\gt{Latent space}{18}{The compressed representation learned in an autoencoder's
bottleneck.}

\gt{Leakage}{7}{Information unavailable at prediction time influencing training.
The most expensive routine error in the field.}

\gt{Learning rate}{4}{Step size in gradient descent. Too large diverges, too small
crawls.}

\gt{Lineage}{23}{The recorded path from source rows to a final number.}

\gt{Logistic regression}{13}{A linear score passed through a sigmoid to give a
probability; coefficients exponentiate to odds ratios.}

\gt{LSTM}{18}{A recurrent cell with gates and an additive cell state, so gradients
survive long sequences.}

\gt{MAE / RMSE}{12}{Mean absolute error treats all errors equally; root mean
squared error punishes large ones.}

\gt{MAR / MCAR / MNAR}{3}{Missingness mechanisms. Under MNAR, encode the absence
rather than imputing it.}

\gt{Membership inference}{26}{Determining whether a record was in a model's
training data --- a privacy attack.}

\gt{Model card}{28}{A short document recording a model's intended use, limitations,
subgroup performance and contest route.}

\gt{Model extraction}{26}{Querying a deployed model repeatedly to train a copy,
then attacking the copy offline.}

\gt{MLOps}{27}{The engineering practice of deploying, monitoring, retraining and
retiring models.}

\gt{Multicollinearity}{8}{Features so correlated they carry the same information;
destabilises linear coefficients.}

\gt{Naive Bayes}{13}{Classification via Bayes' theorem assuming feature
independence. Fast and strong on text.}

\gt{Natural frequencies}{31}{Expressing probability as counts out of a stated
population; understood far better than percentages.}

\gt{Nominal scale}{3}{Unordered categories. Only equality is meaningful.}

\gt{No free lunch}{11}{Averaged over all problems, no algorithm beats any other.}

\gt{Normalisation}{7}{Rescaling to the range zero to one.}

\gt{Odds ratio}{13}{The multiplicative effect of a feature on the odds of the
positive class.}

\gt{One-hot encoding}{7}{Representing a nominal variable as binary columns,
preserving the absence of order.}

\gt{Ordinal scale}{3}{Ordered categories with unequal or unknown gaps.}

\gt{Overfitting}{11}{Learning noise as well as signal; low training error, high
validation error.}

\gt{p-value}{9}{The probability of data this extreme if the null hypothesis were
true. Not the probability the null is true.}

\gt{PCA}{16}{Re-expressing data along orthogonal directions of maximum variance.
Requires scaling; costs interpretability.}

\gt{Pipeline}{7}{A chained sequence of transformations and a model, fitted as one
object so leakage becomes structurally impossible.}

\gt{Poisoning}{26}{Corrupting training data so the next retrained model learns the
attacker's preferred behaviour.}

\gt{Precision}{14}{Of what I flagged, how much was real.}

\gt{Pre-registration}{29}{Recording hypothesis, sample size and analysis plan
before seeing the data.}

\gt{Prompt injection}{22}{Malicious instructions embedded in content an agent
reads. The defining security problem of agentic systems.}

\gt{Provenance}{3}{The documented origin and processing history of a dataset.}

\gt{RAG}{21}{Retrieval-augmented generation: retrieve your documents, put them in
the prompt, cite them in the answer.}

\gt{Random forest}{15}{Bagged decision trees with random feature subsets at each
split. The forgiving default for tabular data.}

\gt{Ratio scale}{3}{Numeric with a true zero; all arithmetic is meaningful.}

\gt{Recall}{14}{Of what was real, how much I caught. Also called sensitivity.}

\gt{Regularisation}{12}{Penalising model complexity to reduce overfitting.}

\gt{ReLU}{17}{The default hidden-layer activation; its gradient does not vanish for
positive inputs.}

\gt{Reproducibility}{29}{Same data and code give the same result. Distinguish from
replicability.}

\gt{Residual}{12}{The gap between actual and predicted. Plot residuals against
fitted values for every regression.}

\gt{Ridge (L2)}{12}{Regularisation that shrinks coefficients without eliminating
them.}

\gt{ROC / AUC}{14}{Recall against false-positive rate across thresholds. Flatters
rare-event detectors --- prefer PR-AUC.}

\gt{Sensor fusion}{25}{Combining several sensors, ideally measuring different
physical quantities, to distinguish sensor faults from process faults.}

\gt{Serverless}{24}{Per-invocation execution with no persistent server. Suits short
stateless work, not training.}

\gt{Shadow deployment}{27}{Running a new model on live traffic without acting on
its output.}

\gt{SHAP}{28}{Per-prediction feature contributions with a consistency guarantee.}

\gt{Sigmoid}{13}{Maps any real number into the interval zero to one.}

\gt{Silhouette score}{16}{Cluster quality measure with a true maximum, so better
than the elbow for choosing the cluster count.}

\gt{Simpson's paradox}{8}{An aggregate relationship that reverses within subgroups.}

\gt{Softmax}{17}{Converts a vector of scores into probabilities summing to one.}

\gt{Spike}{30}{A time-boxed investigation whose deliverable is knowledge.}

\gt{Standardisation}{7}{Rescaling to mean zero and standard deviation one.
Required for distance- and gradient-based methods.}

\gt{Stratified split}{11}{Splitting so class proportions are preserved in every
part. Essential when classes are imbalanced.}

\gt{Supervised learning}{11}{Learning from labelled examples.}

\gt{SVM}{13}{A classifier maximising the margin; kernels allow curved boundaries.}

\gt{Temperature}{21}{Controls sampling sharpness in generation. Near zero for
determinism, higher for variety.}

\gt{TF-IDF}{20}{Scores a term by how distinctive it is to a document relative to
the corpus.}

\gt{Tidy data}{3}{One variable per column, one observation per row, one unit type
per table.}

\gt{Training--serving skew}{27}{Features computed differently in training and
production. Eliminated by a feature store.}

\gt{Transfer learning}{18}{Reusing a model pre-trained on a large dataset and
re-training only its final layers.}

\gt{Transformer}{18}{A stack of self-attention and feed-forward blocks; the basis
of modern language models.}

\gt{Type I / Type II error}{9}{False alarm / missed detection. Which is worse is a
domain question, never a statistical one.}

\gt{UEBA}{26}{User and entity behaviour analytics: scoring each entity against its
own historical baseline.}

\gt{Underfitting}{11}{A model too simple to capture the signal; both errors high.}

\gt{Unsupervised learning}{16}{Finding structure without labels.}

\gt{Validation set}{11}{Data used to choose models and hyperparameters, so the test
set can remain untouched.}

\gt{Variance (ML)}{11}{Error from sensitivity to the particular training sample.}

\gt{Walk-forward validation}{19}{Training on the past and validating on the future.
The only valid scheme for temporal data.}

\gt{Watermark}{19}{The rule for how long a stream waits for late-arriving events
before closing a window.}

\gt{Window (rolling)}{15}{An aggregation over a trailing period. Must never include
future values.}

\end{multicols}
